% Source: https://artofproblemsolving.com/community/c3h1910996_telescoping_questions
% Stable identity: aops-topic-1910996-q01
The $n$-th term is $\dfrac{n+2}{n(n+1)(n+3)(n+4)}$ for $n=1,\dots,10$. Put $m=n+2$, so $m=3,\dots,12$ and the term is
\[
\dfrac{m}{(m-2)(m-1)(m+1)(m+2)}=\dfrac{m}{(m^2-4)(m^2-1)} .
\]
Since $(m^2-1)-(m^2-4)=3$,
\[
\dfrac{m}{(m^2-4)(m^2-1)}=\dfrac{1}{3}\left(\dfrac{m}{m^2-4}-\dfrac{m}{m^2-1}\right),
\]
and $\dfrac{m}{m^2-4}=\dfrac{1}{2}\left(\dfrac{1}{m-2}+\dfrac{1}{m+2}\right)$, $\dfrac{m}{m^2-1}=\dfrac{1}{2}\left(\dfrac{1}{m-1}+\dfrac{1}{m+1}\right)$. Hence the term equals
\[
\dfrac{1}{6}\left[\left(\dfrac{1}{m-2}-\dfrac{1}{m-1}\right)-\left(\dfrac{1}{m+1}-\dfrac{1}{m+2}\right)\right].
\]
Both brackets telescope when summed over $m=3,\dots,12$:
\[
\sum_{m=3}^{12}\left(\dfrac{1}{m-2}-\dfrac{1}{m-1}\right)=1-\dfrac{1}{11}=\dfrac{10}{11},\qquad
\sum_{m=3}^{12}\left(\dfrac{1}{m+1}-\dfrac{1}{m+2}\right)=\dfrac{1}{4}-\dfrac{1}{14}=\dfrac{5}{28}.
\]
So the sum is
\[
\dfrac{1}{6}\left(\dfrac{10}{11}-\dfrac{5}{28}\right)=\dfrac{1}{6}\cdot\dfrac{280-55}{308}=\dfrac{225}{1848}=\boxed{\dfrac{75}{616}} .
\]
